\begin{statementbox}
	\textbf{\textcolor{CStatement}{条件}}
	\begin{description}[style=nextline, leftmargin=2.8em, labelsep=0.5em, itemsep=0.35em]
		\item[\textbf{\textcolor{CStatement}{条件 1（图形类型）：}}]
			$\Omega$ 为中心对称图形（如平行四边形、圆）或三角形 $\triangle ABC$。
	\end{description}
	\textbf{\textcolor{CStatement}{结论}}
	\begin{description}[style=nextline, leftmargin=2.8em, labelsep=0.5em, itemsep=0.45em]
		\item[\textbf{\textcolor{CStatement}{结论一（对称中心平分）：}}]
			\textbf{\textcolor{CStatement}{直观描述：}}
			过对称中心的任意直线平分图形的面积。
			\[
				O\in\ell \;\Rightarrow\; \operatorname{Area}(\Omega\cap H_1)=\operatorname{Area}(\Omega\cap H_2)=\frac12\operatorname{Area}(\Omega).
			\]
		\item[\textbf{\textcolor{CStatement}{结论二（三角形中线平分）：}}]
			\textbf{\textcolor{CStatement}{直观描述：}}
			三角形的中线平分面积，过重心且非中线的直线不一定平分。
			\[
				BM=CM \;\Rightarrow\; \operatorname{Area}(\triangle ABM)=\operatorname{Area}(\triangle ACM)=\frac12\operatorname{Area}(\triangle ABC).
			\]
		\item[\textbf{\textcolor{CStatement}{推广结论（重心判定法）：}}]
			\textbf{\textcolor{CStatement}{直观描述：}}
			过定点 $P$ 和重心 $G$ 的直线，若经过三角形顶点，则平分面积。
			\[
				\exists V\in\{A,B,C\},\; V\in PG\;\Rightarrow\; \operatorname{Area}(\triangle ABC\cap H_1)=\operatorname{Area}(\triangle ABC\cap H_2)=\frac12\operatorname{Area}(\triangle ABC).
			\]
	\end{description}

\end{statementbox}
